\documentclass[../main.tex]{subfiles}

\begin{document}

\chapter{Long Short-Term Memory - LSTM}
\label{chap:LSTM}

\section{Tuning iperparametri su dato pulito}
La colonna timestamp è stata usata per ordinare il dataset in base all'informazione temporale, ma successivamente è stata droppata; questo non comporta una perdita di informazione in quanto rimane intrinseca nell'ordine delle righe.

Gli iperparametri utilizzati per il tuning sono riportati nella Tabella~\ref{tab:lstm_iperparametri}:

\begin{table}[H]
    \centering
    \begin{tabular}{|l|p{5cm}|p{5cm}|}
        \hline
        \textbf{Iperparametro} & \textbf{Descrizione} & \textbf{Valori testati} \\
        \hline
        \texttt{arch\_fc}
          & Architettura degli strati fully-connected dopo le celle LSTM & \texttt{[64]}, \texttt{[128]}, \texttt{[256]}, \texttt{[128, 64]} \\ \hline
        \texttt{lstm\_layers}
          & Numero di strati LSTM impilati & \texttt{1}, \texttt{2} \\ \hline
        \texttt{lr}
          & Learning rate dell'ottimizzatore Adam & \texttt{$10^{-3}$}, \texttt{$10^{-4}$} \\ \hline
        \texttt{epochs}
          & Numero di epoche di addestramento & \texttt{5}, \texttt{10} \\ \hline
    \end{tabular}
    \caption{Iperparametri testati durante il tuning del modello LSTM.}
    \label{tab:lstm_iperparametri}
\end{table}

\paragraph{Scelta dell'ottimizzatore.}
A differenza del modello MLP (Capitolo~\ref{chap:MLP}), che utilizza l'ottimizzatore \textbf{L-BFGS} (\textit{Limited-memory Broyden–Fletcher–Goldfarb–Shanno}) tramite Spark MLlib, il modello LSTM è ottimizzato con \textbf{Adam} (\textit{Adaptive Moment Estimation})~\cite{kingma2017adammethodstochasticoptimization}.

Adam combina il \textit{momentum} del primo momento (media mobile esponenziale dei gradienti) con la normalizzazione del secondo momento (varianza mobile), adattando il learning rate in modo indipendente per ciascun parametro del modello~\cite{kingma2017adammethodstochasticoptimization}. Questa proprietà lo rende particolarmente robusto in presenza di gradienti sparsi o con varianza elevata, condizione frequente nelle celle LSTM che elaborano dipendenze temporali a lungo termine~\cite{ALSELWI2024102068}. Adam è l'ottimizzatore standard per le reti ricorrenti nella letteratura recente ed è il default raccomandato da PyTorch per questo tipo di architettura.

Per il tuning sono stati effettuati tre passaggi importanti riportati di seguito:
\begin{enumerate}
    \item \textbf{Split train/test principale}, effettuato con la stessa logica spiegata nella Sezione~\ref{subsection:datasetSplit};
    \item \textbf{Split tune-train / tune-val};
    \item \textbf{Ordinamento temporale}.
\end{enumerate}

\subsection{Risultati tuning iperparametri}
L'esecuzione del tuning degli iperparametri è durata circa 6 ore per calcolare 30 combinazioni diverse di iperparametri.
I risultati, calcolati con F1-score weighted e AUC-ROC, visibili nella Tabella~\ref{tab:lstm_tuning}, sono in linea con quanto visto nel Capitolo~\ref{chap:MLP} precedente: la F1-Score Weighted rimane pressoché invariata attorno a $\pm 0.96$, con solo alcuni casi in cui scende a $\pm 0.94$ (esecuzioni \textit{22} e \textit{30}). Si è dunque tenuto in considerazione la combinazione di entrambe le metriche e il tempo di esecuzione nel trovare il candidato ideale.

Le tre combinazioni con valore di \textbf{AUC-ROC} più alto sono: combinazione \textbf{6}, \textbf{22} e \textbf{29}. Tuttavia la combinazione \textbf{22} presenta un \textit{F1-score Weighted} di \textit{0.9443}, pertanto è stata scartata. Le due combinazioni rimaste presentano F1 molto simili: rispettivamente $0.9632$ per la combinazione 6 e $0.9624$ per la combinazione 29.
Il fattore decisivo è quindi il tempo: la combinazione 6 esegue in \textbf{687} secondi, mentre la 29 in \textbf{893}.

\begin{table}[H]
\centering
\small
\setlength{\tabcolsep}{5pt}
\begin{tabular}{cccccccc}
\toprule
\textbf{\#} &
\textbf{Arch. FC} &
\textbf{LSTM layers} &
\textbf{LR} &
\textbf{Epoche} &
\textbf{F1} &
\textbf{AUC} &
\textbf{Tempo (s)} \\
\midrule
1  & [64]      & 1 & $10^{-3}$ &  5 & 0.9628 & 0.9409 &   147 \\
2  & [64]      & 1 & $10^{-3}$ & 10 & 0.9628 & 0.9388 &   282 \\
3  & [64]      & 1 & $10^{-4}$ &  5 & 0.9632 & 0.9826 &   143 \\
4  & [64]      & 1 & $10^{-4}$ & 10 & 0.9629 & 0.9787 &   284 \\
5  & [64]      & 2 & $10^{-3}$ &  5 & 0.9628 & 0.9414 &   343 \\
\textbf{6}  & \textbf{[64]} & \textbf{2} & $\mathbf{10^{-3}}$ & \textbf{10} & \textbf{0.9632} & \textbf{0.9933} & \textbf{687} \\
7  & [64]      & 2 & $10^{-4}$ &  5 & 0.9630 & 0.9827 &   341 \\
8  & [64]      & 2 & $10^{-4}$ & 10 & 0.9627 & 0.9467 &   673 \\
\midrule
9  & [128]     & 1 & $10^{-3}$ &  5 & 0.9629 & 0.9360 &   328 \\
10 & [128]     & 1 & $10^{-3}$ & 10 & 0.9629 & 0.9893 &   649 \\
11 & [128]     & 1 & $10^{-4}$ &  5 & 0.9629 & 0.9693 &   328 \\
12 & [128]     & 1 & $10^{-4}$ & 10 & 0.9628 & 0.9519 &   649 \\
13 & [128]     & 2 & $10^{-3}$ &  5 & 0.9631 & 0.9399 &   880 \\
14 & [128]     & 2 & $10^{-3}$ & 10 & 0.9630 & 0.9417 &  1747 \\
15 & [128]     & 2 & $10^{-4}$ &  5 & 0.9632 & 0.9410 &   882 \\
16 & [128]     & 2 & $10^{-4}$ & 10 & 0.9629 & 0.9403 &  1746 \\
\midrule
17 & [256]     & 1 & $10^{-3}$ &  5 & 0.9627 & 0.9481 &   967 \\
18 & [256]     & 1 & $10^{-3}$ & 10 & 0.9630 & 0.9409 &  1915 \\
19 & [256]     & 1 & $10^{-4}$ &  5 & 0.9628 & 0.9438 &   967 \\
20 & [256]     & 1 & $10^{-4}$ & 10 & 0.9630 & 0.9426 &  1916 \\
21 & [256]     & 2 & $10^{-3}$ &  5 & 0.9627 & 0.9594 &  2749 \\
22 & [256]     & 2 & $10^{-3}$ & 10 & 0.9443 & 0.9943 &  5465 \\
23 & [256]     & 2 & $10^{-4}$ &  5 & 0.9631 & 0.9408 &  2756 \\
24 & [256]     & 2 & $10^{-4}$ & 10 & 0.9633 & 0.9406 &  5492 \\
\midrule
25 & [128, 64] & 1 & $10^{-3}$ &  5 & 0.9627 & 0.9630 &   337 \\
26 & [128, 64] & 1 & $10^{-3}$ & 10 & 0.9631 & 0.9695 &   665 \\
27 & [128, 64] & 1 & $10^{-4}$ &  5 & 0.9632 & 0.9652 &   337 \\
28 & [128, 64] & 1 & $10^{-4}$ & 10 & 0.9631 & 0.9464 &   664 \\
29 & [128, 64] & 2 & $10^{-3}$ &  5 & 0.9624 & 0.9915 &   893 \\
30 & [128, 64] & 2 & $10^{-3}$ & 10 & 0.9496 & 0.9391 &  1771 \\
\bottomrule
\end{tabular}
\caption{Risultati del tuning iperparametri LSTM (\texttt{window\_size}~=~30s).
La metrica di selezione è AUC-ROC: l'F1 weighted risulta non discriminante per via
dello sbilanciamento delle classi ($\sim$1.3\% guasti), mentre l'AUC varia
significativamente tra le configurazioni (0.936--0.994).
La configurazione ottimale (\#6, evidenziata) è \texttt{arch=[64]},
\texttt{lstm\_layers=2}, \texttt{lr=$10^{-3}$}, \texttt{epochs=10},
con AUC=0.9933.}
\label{tab:lstm_tuning}
\end{table}

Questa scelta è stata anche convalidata visivamente tramite due grafici.

La Figura~\ref{fig:LSTMefficienzaTuning} mostra il rapporto tra la metrica \textbf{AUC-ROC} e il tempo di esecuzione in secondi. Si può notare come il valore AUC migliore appartenga alla combinazione \textbf{22}, precedentemente scartata per via dell'F1 insufficiente. Questa combinazione è anche l'unica ad aver impiegato più di 5000 secondi ($\approx 83$ minuti), rispetto ai $\approx 11$ minuti della combinazione \textbf{6}.

\begin{figure}[H]
    \centering
    \includegraphics[width=1\linewidth]{images/Results/lstm_tuning_efficiency.png}
    \caption{Confronto tra AUC-ROC e tempo di esecuzione per il tuning iperparametri LSTM.}
    \label{fig:LSTMefficienzaTuning}
\end{figure}

La Figura~\ref{fig:LSTMConfrontoTuningMetriche} mostra invece il confronto visivo tra le due metriche, confermando quanto detto sopra: la combinazione \textbf{22}, pur essendo quella con AUC-ROC migliore, risulta anche quella con F1 peggiore.

\begin{figure}[H]
    \centering
    \includegraphics[width=1\linewidth]{images/Results/lstm_tuning_metrics.png}
    \caption{Confronto tra AUC-ROC e F1-Score Weighted per il tuning iperparametri LSTM.}
    \label{fig:LSTMConfrontoTuningMetriche}
\end{figure}

% -----------------------------------------------------------------------
\section{Esecuzione del modello}
\label{section:LSTM}

A differenza dell'MLP, il training dell'LSTM è stato eseguito \textbf{localmente} su GPU dedicata (NVIDIA GTX 1070), anziché sul cluster Apache Spark. La scelta è motivata da due ragioni principali.

Dal punto di vista architetturale, l'LSTM non è integrato nativamente in MLlib: l'unica via di integrazione con Spark è tramite \texttt{TorchDistributor}, un wrapper che coordina job PyTorch distribuiti sui worker del cluster. Tuttavia, il deployment di \texttt{TorchDistributor} richiede che ogni worker disponga delle stesse dipendenze Python (PyTorch, CUDA toolkit, librerie ausiliarie), condizione non soddisfatta nell'ambiente di cluster utilizzato, dove i tentativi di installazione hanno esaurito lo spazio su disco prima del completamento.

Dal punto di vista delle prestazioni, anche senza i problemi di configurazione, l'architettura LSTM con finestre temporali di 30 secondi risulta intrinsecamente sequenziale: il calcolo degli stati nascosti non è parallelizzabile lungo la dimensione temporale, rendendo il vantaggio del cluster distribuito marginale rispetto a una singola GPU locale. L'esecuzione su GTX 1070 ha impiegato circa \textbf{113,2 secondi per run}, consentendo di completare le 1100 run totali in circa \textbf{35 ore} di calcolo cumulativo.\footnote{Il tempo indicato si riferisce alla somma dei tempi di esecuzione sequenziale. È stato tentato un approccio parallelo lanciando due istanze simultanee, ma con soli 16 GB di RAM disponibili il sistema andava frequentemente in paging o terminava con errori di memoria, rendendo le esecuzioni inaffidabili. Si è quindi optato per l'esecuzione sequenziale, che garantisce stabilità a scapito del tempo totale.}

% -----------------------------------------------------------------------
\section{Risultati}

La baseline del modello LSTM, calcolata su 20 esecuzioni indipendenti senza rumore (Tabella~\ref{tab:lstm_baseline}), è $\text{F1} = 0.9739 \pm 0.0050$ e $\text{AUC} = 0.9638 \pm 0.0101$. Rispetto alla baseline MLP ($\text{AUC} = 0.9591$, Tabella~\ref{tab:mlp_baseline}), l'LSTM parte da un punto di partenza comparabile in termini di capacità discriminante: sfruttando l'informazione temporale codificata nelle finestre di 30 secondi, il modello è in grado di rilevare le anomalie di pressione con maggiore sensibilità già sul dato pulito.

Il comportamento sotto rumore replica in parte le tendenze già osservate per MLP nel Capitolo~\ref{chap:MLP}, con alcune differenze strutturali rilevanti. Come per MLP, il rumore \textit{labels} produce la degradazione più severa, mentre \textit{duplicated} non produce alcun effetto misurabile. I metodi \textit{missing}, \textit{noise} e \textit{outliers} applicati alle feature mostrano invece un comportamento parzialmente diverso rispetto a MLP, in particolare per quanto riguarda le feature \textbf{TP3} e \textbf{Oil\_temperature}, come analizzato nelle sezioni successive.

Tutti i risultati sono riportati nelle Tabelle~\ref{tab:lstm_baseline}, \ref{tab:lstm_duplicated}, \ref{tab:lstm_labels}, \ref{tab:lstm_missing}, \ref{tab:lstm_noise} e~\ref{tab:lstm_outliers} e nelle Figure~\ref{fig:LSTMf1Degradazione} e~\ref{fig:LSTMAUCDegradazione}.

\newpage

\subsection{Labels}
Come mostrato nella Sezione~\ref{subsec:corruzione_labels}, al 50\% di corruzione le due classi si equivalgono esattamente, privando anche l'LSTM di qualsiasi segnale discriminativo affidabile. Il rumore \textit{labels} produce, anche sull'LSTM, la degradazione più severa tra tutte le tipologie testate, in linea con quanto documentato da Zhang et al.~\cite{zhang2017understandingdeeplearningrequires}. Tuttavia, il collasso è ancora più pronunciato rispetto a MLP: a 50\% di rumore, l'AUC-ROC crolla da una baseline di $0.9638$ fino a $\mathbf{0.3511 \pm 0.0981}$, un valore inferiore persino a quello di un classificatore casuale ($0.5$), indicando che il modello ha invertito il proprio ordinamento interno delle classi. A titolo comparativo, l'MLP alla stessa percentuale si attesta a $0.5119 \pm 0.1889$ (Tabella~\ref{tab:mlp_labels}).

L'F1-score weighted subisce invece il crollo con un ritardo di una soglia rispetto all'AUC: rimane pressoché stabile fino al 30\% ($0.9798 \pm 0.0086$) per poi collassare a $0.4681 \pm 0.1629$ al 50\%. Questo disallineamento temporale tra le due metriche evidenzia che l'LSTM, sfruttando il contesto temporale delle finestre, riesce a mantenere un minimo di coerenza nell'assegnazione della classe maggioritaria più a lungo di MLP, ma una volta superata la soglia critica degrada in modo più brusco.

\subsection{Duplicated}
Come analizzato nella Sezione~\ref{subsec:corruzione_duplicated}, la duplicazione casuale non altera la distribuzione delle classi e lo sbilanciamento 98.7\%/1.3\% rimane invariato a qualsiasi percentuale. Come per MLP, il modello LSTM risulta completamente resistente all'introduzione di righe duplicate. La F1-score rimane stabile o leggermente superiore alla baseline entro l'intervallo di confidenza, con la massima deviazione di $+0.0060$ a 10\% di rumore. Anche l'AUC-ROC segue lo stesso andamento, con valori compresi nell'intervallo $[0.9486, 0.9735]$, tutti sovrapposti alla baseline ($0.9638 \pm 0.0101$).

La resistenza ai duplicati è spiegabile dalla natura del metodo: le righe aggiunte non alterano la distribuzione informativa del dataset ma ne aumentano la ridondanza. Per un modello sequenziale come LSTM, che elabora finestre di campioni, l'introduzione di finestre duplicate produce effettivamente un overfit locale su quei pattern, ma la quantità di pattern distinti rimasta è sufficiente a preservare la capacità discriminante globale.

% -----------------------------------------------------------------------
\section{F1-Score}

Come già osservato per MLP nel Capitolo~\ref{chap:MLP}, l'F1-score weighted risulta una metrica scarsamente sensibile alla corruzione delle feature su questo dataset, a causa dello sbilanciamento delle classi~\cite{buda2018systematic, johnson2019survey}. Il 98\% dei campioni appartiene alla classe \textit{senza guasti} e il contributo ponderato di questa classe domina la metrica aggregata, mascherando degradazioni anche significative sulla classe minoritaria.

Nei metodi \textit{missing}, \textit{noise} e \textit{outliers}, la F1-score rimane sostanzialmente stabile per tutte e tre le feature testate: la deviazione massima dalla baseline a 50\% di rumore è di $0.0149$ punti, registrata per \textit{missing} su \texttt{TP3} ($0.9590 \pm 0.0041$, Tabella~\ref{tab:lstm_missing}). Gli intervalli di confidenza al 95\% si sovrappongono quasi completamente a quelli della condizione pulita, indicando assenza di degradazione statisticamente significativa sull'F1.

L'unica eccezione è il rumore \texttt{Labels}, dove entrambe le metriche degradano in modo netto e statisticamente significativo, come visibile nella Figura~\ref{fig:LSTMf1Degradazione}. Il crollo dell'F1 da $0.9739 \pm 0.0050$ a $0.4681 \pm 0.1629$ al 50\% di rumore (Tabella~\ref{tab:lstm_labels}) è una conseguenza diretta dell'alterazione del segnale di supervisione durante il training, confermando quanto già documentato per MLP e per i modelli deep in generale~\cite{zhang2017understandingdeeplearningrequires}.

% -----------------------------------------------------------------------
\section{AUC-ROC}

L'analisi dell'AUC-ROC per i metodi \textit{missing}, \textit{noise} e \textit{outliers} rivela un comportamento parzialmente diverso rispetto a MLP, con differenze rilevanti nelle feature \textbf{TP3} e \textbf{Oil\_temperature}.

\subsection{DV\_pressure}
Il comportamento di \texttt{DV\_pressure} sull'LSTM è analogo a quanto osservato per MLP: la feature presenta una distribuzione quasi-binaria (98\% dei valori a 0, con 5 durante i guasti), e la sua struttura informativamente netta la rende resistente a tutti i tipi di corruzione. L'AUC-ROC rimane stabile o migliora per \textit{noise} ($0.9788 \pm 0.0089$ a 50\%) e per \textit{outliers}. Per \textit{missing} si osserva una leggera oscillazione, ma i valori rimangono all'interno degli intervalli di confidenza della baseline.

La spiegazione è la stessa già data per MLP: il metodo \textit{noise} genera valori uniformi in $[0, 5]$, producendo valori medi attorno a $2.5$, sufficientemente distanti dal valore di guasto ($5$) da non compromettere la separabilità delle classi. La struttura temporale delle finestre LSTM non altera questa proprietà intrinseca della feature. È inoltre noto che l'aggiunta di rumore ai dati di input durante il training agisce come forma implicita di regolarizzazione~\cite{bishop1995training}, il che può spiegare la stabilità o il lieve miglioramento dell'AUC osservato su questa feature altamente strutturata.

\subsection{TP3}
La feature \texttt{TP3} mostra sull'LSTM un comportamento diverso rispetto a MLP. Come illustrato nel Capitolo~\ref{chap:MLP}, l'MLP risultava quasi immune alla corruzione di \texttt{TP3} grazie alla correlazione perfetta con \texttt{Reservoirs}: quando \texttt{TP3} veniva corrotta, il modello si appoggiava automaticamente a \texttt{Reservoirs} come feature sostitutiva.

Sull'LSTM questo meccanismo di fallback risulta meno efficace: il modello elabora finestre temporali dove la struttura sequenziale è parte integrante del segnale, e la sostituzione implicita tra feature correlate è meno immediata. Il risultato è una degradazione misurabile dell'AUC: con \textit{missing} al 50\% l'AUC scende a $0.9421 \pm 0.0237$ (vs $0.9524 \pm 0.0102$ di MLP, Tabella~\ref{tab:mlp_missing}), con intervalli di confidenza che si sovrappongono appena alla baseline ($0.9638 \pm 0.0101$).

Tuttavia, la degradazione rimane contenuta rispetto a \texttt{Oil\_temperature}, suggerendo che la correlazione con \texttt{Reservoirs} attenua comunque parzialmente l'impatto, pur senza eliminarlo completamente come accadeva per MLP.

\subsection{Oil\_temperature}
Il comportamento di \texttt{Oil\_temperature} sull'LSTM è, sorprendentemente, più robusto rispetto a MLP. Mentre su MLP questa feature produceva la degradazione più severa tra le tre (AUC fino a $0.8137 \pm 0.0487$ con \textit{missing} al 50\%, Tabella~\ref{tab:mlp_missing}), sull'LSTM i valori rimangono in un range molto più alto, con il minimo registrato a $0.9460 \pm 0.0169$ (\textit{missing} 50\%) e $0.9224 \pm 0.0227$ (\textit{outliers} 30\%).

La spiegazione più probabile risiede nella capacità dell'LSTM di sfruttare il contesto temporale per compensare la corruzione di singoli valori. \texttt{Oil\_temperature} presenta una dinamica lenta e continua nel tempo (cresce gradualmente durante le fasi di carico, vedi Capitolo~\ref{chap:DatasetAnalysis}): il modello, elaborando finestre di 30 secondi, può ricostruire implicitamente il valore atteso dai campioni adiacenti nella finestra stessa, attenuando l'impatto dei valori corrotti o mancanti. L'MLP, non disponendo di questo contesto temporale, dipende esclusivamente dal valore istantaneo della feature e risulta quindi più vulnerabile.

\begin{figure}[H]
    \centering
    \includegraphics[width=1\linewidth]{images/Results/lstm_f1_degradation.png}
    \caption{Degradazione del modello LSTM al crescere della percentuale di rumore, metrica F1-score.}
    \label{fig:LSTMf1Degradazione}
\end{figure}

\begin{figure}[H]
    \centering
    \includegraphics[width=1\linewidth]{images/Results/lstm_auc_degradation.png}
    \caption{Degradazione del modello LSTM al crescere della percentuale di rumore, metrica AUC-ROC.}
    \label{fig:LSTMAUCDegradazione}
\end{figure}

\newpage



\section{Esperimenti di approfondimento}
\label{sec:lstm_approfondimento}

\subsection{Verifica della ridondanza strutturale: TP3 e Reservoir}
\label{subsec:lstm_tp3_reservoir}

L'esperimento è stato progettato per verificare se la ridondanza strutturale tra \texttt{TP3} e \texttt{Reservoirs}, già osservata per MLP nella Sezione~\ref{subsec:tp3_reservoir}, operi in modo analogo sull'LSTM. Nell'esperimento principale, la corruzione veniva applicata a una sola feature alla volta: in quella condizione, il modello poteva appoggiarsi implicitamente alla feature correlata come segnale sostitutivo. Applicando il rumore \textit{simultaneamente} a entrambe le feature, si elimina questo meccanismo di fallback, isolando la dipendenza del modello dall'informazione condivisa dalle due feature.

I risultati sono riportati nelle Tabelle~\ref{tab:lstm_new_duplicated_minority} e~\ref{tab:lstm_new_tp3_reservoir} e nelle Figure~\ref{fig:LSTMNewF1Degradazione} e~\ref{fig:LSTMNewAUCDegradazione}.

\paragraph{Missing.}
La corruzione congiunta produce una degradazione progressiva e sistematica dell'AUC-ROC, che scende da $0.9638$ (baseline) fino a $0.8592 \pm 0.0590$ al 50\% di rumore. Il confronto con l'esperimento su \texttt{TP3} singola (Tabella~\ref{tab:lstm_missing}), dove l'AUC al 50\% si attestava a $0.9421 \pm 0.0237$, conferma che la degradazione aggiuntiva di $\approx 0.08$ punti è direttamente attribuibile alla rimozione del meccanismo di fallback su \texttt{Reservoirs}: privato di entrambe le sorgenti informative, il modello non riesce a compensare i valori mancanti nemmeno sfruttando il contesto temporale delle finestre. %La stessa tendenza è osservabile anche per MLP (Tabella~\ref{tab:mlp_tp3_reservoir}), dove l'AUC al 50\% con \textit{missing} si attesta a $0.9086 \pm 0.0224$: in quel caso però la degradazione è più contenuta, a conferma che l'LSTM — pur essendo più robusto sulla feature singola grazie al contesto temporale — risulta qui più vulnerabile, poiché il segnale sequenziale che sfruttava come compensazione si basa anch'esso su \texttt{TP3} e \texttt{Reservoirs}.

\paragraph{Noise.}
Il rumore \textit{noise} produce il risultato più rilevante dell'intero esperimento. L'AUC-ROC \textit{migliora} rispetto alla baseline nelle prime tre soglie: $0.9693 \pm 0.0180$ al 10\%, $0.9903 \pm 0.0049$ al 20\%, $0.9906 \pm 0.0045$ al 30\%, per poi collassare bruscamente a $0.7629 \pm 0.0509$ al 50\%. Anche l'F1-score, rimasta stabile fino al 30\%, subisce un crollo netto a $0.8914 \pm 0.0384$.

Il miglioramento osservato alle basse percentuali è interpretabile come un effetto di \textbf{regolarizzazione implicita}: l'aggiunta di rumore ai valori di input durante il training introduce una perturbazione che contrasta l'overfitting, migliorando la capacità di generalizzazione del modello~\cite{10.1162/neco.1995.7.1.108}. Questo effetto benefico è tuttavia limitato a intensità di rumore moderate; quando la percentuale raggiunge il 50\%, la corruzione diventa sufficientemente severa da distruggere la struttura temporale delle sequenze su cui l'LSTM si basa, e le performance collassano in modo brusco e non recuperabile.% Il confronto con MLP (Tabella~\ref{tab:mlp_tp3_reservoir}), dove il metodo \textit{noise} produce invece una degradazione monotona e graduale dell'AUC, evidenzia come la sensibilità dell'LSTM alla coerenza temporale del segnale amplifichi sia il beneficio della regolarizzazione a bassa intensità, sia il crollo a intensità elevata.

\paragraph{Outliers.}
Il comportamento con \textit{outliers} è non monotono: l'AUC scende a $0.8864 \pm 0.0325$ al 10\% per poi oscillare attorno a $0.88$--$0.90$ fino al 30\%, e recuperare a $0.9309 \pm 0.0102$ al 50\%. L'F1-score rimane sostanzialmente stabile per tutte le percentuali, con variazioni inferiori al margine di confidenza. Il recupero dell'AUC al 50\% suggerisce che, ad alta densità di outlier, il modello ridistribuisce implicitamente i pesi su altre feature non corrotte, attenuando l'impatto della corruzione congiunta. Questo comportamento è analogo a quanto osservato per MLP nella stessa condizione (Tabella~\ref{tab:mlp_tp3_reservoir}).

\subsection{Duplicated sulla classe minoritaria}
\label{subsec:lstm_duplicated_minority}

Come documentato nella Sezione~\ref{subsec:corruzione_duplicated}, anche con oversampling al 50\% sulla classe minoritaria la quota di guasti passa solo dall'$1.29\%$ all'$1.92\%$, un incremento troppo contenuto per alterare il bilanciamento globale del dataset.

I risultati (Tabella~\ref{tab:lstm_new_duplicated_minority}) mostrano un comportamento diverso rispetto all'esperimento \textit{duplicated} standard (Tabella~\ref{tab:lstm_duplicated}), dove sia F1 che AUC rimanevano stabili entro l'intervallo di confidenza della baseline. Qui invece l'AUC-ROC degrada al crescere della percentuale di rumore: da $0.9638$ (baseline) a $0.9539 \pm 0.0166$ al 10\%, fino a $0.9166 \pm 0.0391$ al 30\%, per poi stabilizzarsi a $0.9196 \pm 0.0354$ al 50\%. L'F1-score rimane invece pressoché invariata, confermando il comportamento già osservato per MLP~\cite{buda2018systematic}.
La degradazione dell'AUC è spiegabile con il fatto che la duplicazione mirata sulla classe minoritaria (classe guasti) introduce sequenze ripetute di eventi $target = 1$ nelle finestre temporali elaborate dall'LSTM. Queste sequenze duplicate incrementano artificialmente il peso dei pattern di guasto durante il training, senza però migliorare il bilanciamento reale del dataset che rimane quasi invariato. Il risultato è un overfitting selettivo sulla classe minoritaria che degrada il modello, riducendone l'AUC senza impattare l'F1, dominata dalla classe maggioritaria~\cite{johnson2019survey}.

%A differenza dell'MLP (Tabella~\ref{tab:mlp_duplicated_minority}), dove la stessa procedura produceva una degradazione dell'AUC molto più contenuta ($0.9514 \pm 0.0093$ al 10\%, $0.9466 \pm 0.0079$ al 50\%), sull'LSTM la riduzione è più marcata, in particolare al 30\% ($0.9166 \pm 0.0391$). Questo suggerisce che la struttura sequenziale delle finestre amplifichi l'effetto della ridondanza locale: un campione duplicato nell'MLP contribuisce come singolo vettore di feature; nell'LSTM lo stesso campione può comparire in più finestre consecutive, moltiplicando la sua influenza sul gradiente durante il backpropagation through time~\cite{Hochreiter1997}.


\begin{figure}[H]
    \centering
    \includegraphics[width=1\linewidth]{images/Results/lstm_new_f1_degradation.png}
    \caption{Degradazione dell'F1-score del modello LSTM al crescere della percentuale di rumore, esperimento con feature aggregate \texttt{TP3 + Reservoirs} e \texttt{Target (fault=1)}. 20 run per punto, IC 95\%.}
    \label{fig:LSTMNewF1Degradazione}
\end{figure}

\begin{figure}[H]
    \centering
    \includegraphics[width=1\linewidth]{images/Results/lstm_new_auc_degradation.png}
    \caption{Degradazione dell'AUC-ROC del modello LSTM al crescere della percentuale di rumore, esperimento con feature aggregate \texttt{TP3 + Reservoirs} e \texttt{Target (fault=1)}. 20 run per punto, IC 95\%.}
    \label{fig:LSTMNewAUCDegradazione}
\end{figure}


\newpage


% -----------------------------------------------------------------------
\section{Tabelle di riferimento per LSTM}
\begin{table}[H]
\centering\small
\rowcolors{2}{gray!10}{white}
\begin{tabular}{ccc}
\toprule
\textbf{Noise \%} & \textbf{F1 (media $\pm$ IC 95\%)} & \textbf{AUC (media $\pm$ IC 95\%)} \\
\midrule
0\%  & $0.9739 \pm 0.0050$ & $0.9638 \pm 0.0101$ \\
\bottomrule
\end{tabular}
\caption{Tabella risultati LSTM, baseline del modello calcolata su 20 esecuzioni senza rumore.}
\label{tab:lstm_baseline}
\end{table}
%--------------------------------
\begin{table}[H]
\centering\small
\rowcolors{2}{gray!10}{white}
\begin{tabular}{cccc}
\toprule
\textbf{Noise \%} & \textbf{Feature} & \textbf{F1 (media $\pm$ IC 95\%)} & \textbf{AUC (media $\pm$ IC 95\%)} \\
\midrule
10\% & Row based & $0.9799 \pm 0.0037$ & $0.9688 \pm 0.0124$ \\
20\% & Row based & $0.9769 \pm 0.0034$ & $0.9486 \pm 0.0240$ \\
30\% & Row based & $0.9777 \pm 0.0031$ & $0.9735 \pm 0.0107$ \\
50\% & Row based & $0.9775 \pm 0.0039$ & $0.9666 \pm 0.0114$ \\
\bottomrule
\end{tabular}
\caption{Tabella risultati LSTM, rumore dal metodo \textit{duplicated}.}
\label{tab:lstm_duplicated}
\end{table}
%--------------------------------
\begin{table}[H]
\centering\small
\rowcolors{2}{gray!10}{white}
\begin{tabular}{cccc}
\toprule
\textbf{Noise \%} & \textbf{Feature} & \textbf{F1 (media $\pm$ IC 95\%)} & \textbf{AUC (media $\pm$ IC 95\%)} \\
\midrule
10\% & Target & $0.9793 \pm 0.0044$ & $0.6700 \pm 0.1233$ \\
20\% & Target & $0.9789 \pm 0.0057$ & $0.7563 \pm 0.1119$ \\
30\% & Target & $0.9798 \pm 0.0086$ & $0.6876 \pm 0.1046$ \\
50\% & Target & $\mathbf{0.4681 \pm 0.1629}$ & $\mathbf{0.3511 \pm 0.0981}$ \\
\bottomrule
\end{tabular}
\caption{Tabella risultati LSTM, rumore dal metodo \textit{labels}.}
\label{tab:lstm_labels}
\end{table}
%--------------------------------
\begin{table}[H]
\centering\small
\rowcolors{2}{gray!10}{white}
\begin{tabular}{cccc}
\toprule
\textbf{Noise \%} & \textbf{Feature} & \textbf{F1 (media $\pm$ IC 95\%)} & \textbf{AUC (media $\pm$ IC 95\%)} \\
\midrule
10\% & DV\_pressure     & $0.9630 \pm 0.0052$ & $0.9651 \pm 0.0153$ \\
20\% & DV\_pressure     & $0.9675 \pm 0.0049$ & $0.9855 \pm 0.0041$ \\
30\% & DV\_pressure     & $0.9611 \pm 0.0035$ & $0.9582 \pm 0.0200$ \\
50\% & DV\_pressure     & $0.9721 \pm 0.0058$ & $0.9779 \pm 0.0065$ \\
\midrule
10\% & Oil\_temperature & $0.9764 \pm 0.0025$ & $0.9485 \pm 0.0201$ \\
20\% & Oil\_temperature & $0.9809 \pm 0.0025$ & $0.9668 \pm 0.0094$ \\
30\% & Oil\_temperature & $0.9773 \pm 0.0040$ & $0.9508 \pm 0.0267$ \\
50\% & Oil\_temperature & $0.9793 \pm 0.0020$ & $0.9460 \pm 0.0169$ \\
\midrule
10\% & TP3             & $0.9695 \pm 0.0057$ & $0.9413 \pm 0.0278$ \\
20\% & TP3             & $0.9600 \pm 0.0040$ & $0.9428 \pm 0.0289$ \\
30\% & TP3             & $0.9590 \pm 0.0040$ & $0.9524 \pm 0.0232$ \\
50\% & TP3             & $0.9590 \pm 0.0041$ & $0.9421 \pm 0.0237$ \\
\bottomrule
\end{tabular}
\caption{Tabella risultati LSTM, rumore dal metodo \textit{missing}.}
\label{tab:lstm_missing}
\end{table}
%--------------------------------
\begin{table}[H]
\centering\small
\rowcolors{2}{gray!10}{white}
\begin{tabular}{cccc}
\toprule
\textbf{Noise \%} & \textbf{Feature} & \textbf{F1 (media $\pm$ IC 95\%)} & \textbf{AUC (media $\pm$ IC 95\%)} \\
\midrule
10\% & DV\_pressure     & $0.9728 \pm 0.0050$ & $0.9750 \pm 0.0096$ \\
20\% & DV\_pressure     & $0.9783 \pm 0.0036$ & $0.9750 \pm 0.0123$ \\
30\% & DV\_pressure     & $0.9766 \pm 0.0056$ & $0.9811 \pm 0.0063$ \\
50\% & DV\_pressure     & $0.9822 \pm 0.0030$ & $0.9788 \pm 0.0089$ \\
\midrule
10\% & Oil\_temperature & $0.9746 \pm 0.0044$ & $0.9373 \pm 0.0207$ \\
20\% & Oil\_temperature & $0.9734 \pm 0.0051$ & $0.9668 \pm 0.0129$ \\
30\% & Oil\_temperature & $0.9763 \pm 0.0045$ & $0.9324 \pm 0.0243$ \\
50\% & Oil\_temperature & $0.9791 \pm 0.0034$ & $0.9547 \pm 0.0178$ \\
\midrule
10\% & TP3             & $0.9689 \pm 0.0058$ & $0.9754 \pm 0.0111$ \\
20\% & TP3             & $0.9754 \pm 0.0043$ & $0.9746 \pm 0.0079$ \\
30\% & TP3             & $0.9786 \pm 0.0048$ & $0.9655 \pm 0.0134$ \\
50\% & TP3             & $0.9732 \pm 0.0055$ & $0.9822 \pm 0.0057$ \\
\bottomrule
\end{tabular}
\caption{Tabella risultati LSTM, rumore dal metodo \textit{noise}.}
\label{tab:lstm_noise}
\end{table}
%--------------------------------
\begin{table}[H]
\centering\small
\rowcolors{2}{gray!10}{white}
\begin{tabular}{cccc}
\toprule
\textbf{Noise \%} & \textbf{Feature} & \textbf{F1 (media $\pm$ IC 95\%)} & \textbf{AUC (media $\pm$ IC 95\%)} \\
\midrule
10\% & DV\_pressure     & $0.9732 \pm 0.0041$ & $0.9535 \pm 0.0221$ \\
20\% & DV\_pressure     & $0.9779 \pm 0.0039$ & $0.9787 \pm 0.0068$ \\
30\% & DV\_pressure     & $0.9765 \pm 0.0046$ & $0.9767 \pm 0.0135$ \\
50\% & DV\_pressure     & $0.9705 \pm 0.0062$ & $0.9817 \pm 0.0112$ \\
\midrule
10\% & Oil\_temperature & $0.9749 \pm 0.0049$ & $0.9539 \pm 0.0185$ \\
20\% & Oil\_temperature & $0.9767 \pm 0.0046$ & $0.9727 \pm 0.0090$ \\
30\% & Oil\_temperature & $0.9690 \pm 0.0045$ & $0.9224 \pm 0.0227$ \\
50\% & Oil\_temperature & $0.9627 \pm 0.0047$ & $0.9032 \pm 0.0692$ \\
\midrule
10\% & TP3             & $0.9750 \pm 0.0049$ & $0.9567 \pm 0.0188$ \\
20\% & TP3             & $0.9709 \pm 0.0052$ & $0.9691 \pm 0.0068$ \\
30\% & TP3             & $0.9703 \pm 0.0059$ & $0.9653 \pm 0.0131$ \\
50\% & TP3             & $0.9630 \pm 0.0050$ & $0.9720 \pm 0.0117$ \\
\bottomrule
\end{tabular}
\caption{Tabella risultati LSTM, rumore dal metodo \textit{outliers}.}
\label{tab:lstm_outliers}
\end{table}
%--------------------------------
\begin{table}[H]
\centering\small
\rowcolors{2}{gray!10}{white}
\begin{tabular}{cccc}
\toprule
\textbf{Noise \%} & \textbf{Feature} & \textbf{F1 (media $\pm$ IC 95\%)} & \textbf{AUC (media $\pm$ IC 95\%)} \\
\midrule
10\% & Target (1) & $0.9749 \pm 0.0038$ & $0.9539 \pm 0.0166$ \\
20\% & Target (1) & $0.9722 \pm 0.0045$ & $0.9421 \pm 0.0157$ \\
30\% & Target (1) & $0.9711 \pm 0.0045$ & $0.9166 \pm 0.0391$ \\
50\% & Target (1) & $0.9716 \pm 0.0054$ & $0.9196 \pm 0.0354$ \\
\bottomrule
\end{tabular}
\caption{Tabella risultati LSTM, rumore \textit{duplicated} applicato
esclusivamente alla classe minoritaria (target = 1).
La baseline di riferimento è riportata in Tabella~\ref{tab:lstm_baseline}.}
\label{tab:lstm_new_duplicated_minority}
\end{table}

\begin{table}[H]
\centering\small
\rowcolors{2}{gray!10}{white}
\begin{tabular}{cccc}
\toprule
\textbf{Noise \%} & \textbf{Metodo} & \textbf{F1 (media $\pm$ IC 95\%)} & \textbf{AUC (media $\pm$ IC 95\%)} \\
\midrule
10\% & missing  & $0.9546 \pm 0.0026$ & $0.9289 \pm 0.0191$ \\
20\% & missing  & $0.9531 \pm 0.0047$ & $0.9138 \pm 0.0256$ \\
30\% & missing  & $0.9497 \pm 0.0085$ & $0.8767 \pm 0.0392$ \\
50\% & missing  & $0.9521 \pm 0.0053$ & $0.8592 \pm 0.0590$ \\
\midrule
10\% & noise    & $0.9800 \pm 0.0039$ & $0.9693 \pm 0.0180$ \\
20\% & noise    & $0.9825 \pm 0.0017$ & $0.9903 \pm 0.0049$ \\
30\% & noise    & $0.9729 \pm 0.0049$ & $0.9906 \pm 0.0045$ \\
50\% & noise    & $0.8914 \pm 0.0384$ & $0.7629 \pm 0.0509$ \\
\midrule
10\% & outliers & $0.9550 \pm 0.0026$ & $0.8864 \pm 0.0325$ \\
20\% & outliers & $0.9573 \pm 0.0055$ & $0.8973 \pm 0.0168$ \\
30\% & outliers & $0.9566 \pm 0.0054$ & $0.8810 \pm 0.0377$ \\
50\% & outliers & $0.9601 \pm 0.0042$ & $0.9309 \pm 0.0102$ \\
\bottomrule
\end{tabular}
\caption{Tabella risultati LSTM, esperimento con rumore applicato
simultaneamente su \texttt{TP3} e \texttt{Reservoirs}.
La baseline di riferimento è riportata in Tabella~\ref{tab:lstm_baseline}.}
\label{tab:lstm_new_tp3_reservoir}
\end{table}
%--------------------------------
\newpage

\end{document}